\section{What a finite lottery certificate buys in Chile}\label{sec:lottery-tests}
This appendix reports the original pooled-seat benchmark. Appendix~\ref{sec:institution} gives the corresponding institutional results. The common options and entitlements on the two frozen menus stay fixed. Suppose there are $d_i$ old-only and $d_i$ proposed-only options for pupil $i$. Among the old-only options let $p_i$ count those that the pupil receives in any protected draw. Among the proposed-only options let $f_i$ count those blocking at least one protected allocation. At most $d_i-p_i$ old slots can be freed and at most $d_i-f_i$ new slots can be used. Any eligible new slot can replace any removable old one for this pupil. Thus the exact number of substitutions toward the new policy is $d_i-\max(p_i,f_i)$. We use submitted rank to break ties. The objective counts proposed choices adopted. An additional choice need not be a useful prediction.

\paragraph{Design}
For each of 2022--24 we drew 16 priority profiles using independent within-course permutations inside the five public priority categories. We solved the zero-, one- and two-exception problems exactly, enumerating the possible draws to exclude. We evaluated on the recorded within-course order and on 64 fresh draws shared by baseline and challengers. We fixed these choices before computing the new outcomes on these previously studied cohorts. With exception budgets of zero, one and two, the optimal menus retained 39,293, 39,727 and 40,189 proposed substitutions, respectively. Under the recorded orders 146, 158 and 184 former recipients lost their places. The entire switch made 67,909 substitutions and displaced 2,088. A simulation batch is a specification of possible priorities rather than an estimate of Chile's actual priority law. We exclude 2025 because the tie procedure changed.

\paragraph{Training support and edit-matched comparisons}
After seeing the first test we extended the nested training support to 64, 256 and 1,024 draws, always requiring the old assignment to be stable on all training draws. For each resulting menu we made an alternative with exactly the same number of slot substitutions and changed rows. We permuted the numbers of swaps among pupils with the same number of differing slots, then sampled which old-only and new-only options to exchange. This control retains all shared courses and protected options. Eight seeds are evaluated at the recorded priorities, and the first seed also on every simulated test draw.

\begin{table}[ht]\centering\small
\caption{Nested training supports on the same 64 test draws per year. Failure counts pool 192 cohort--lottery realizations. A certificate fails when a baseline pair is missing or has a blocking pair. A gross exit is a baseline recipient who loses their placement. Repeat appearances in different lotteries count separately.}\label{tab:nested-risk}
\begin{tabular}{@{}rrrrrr@{}}\toprule
Design draws & Adopted slots & Failed certificates & Draws with exits & Gross exits & Control exits\\\midrule
16 &39,293&192&192&9,491&78,315\\
64 &33,452&192&189&1,087&67,499\\
256 &30,709&161&103&161&62,575\\
1,024 &29,018&52&30&33&59,768\\\bottomrule
\end{tabular}
\end{table}

On the recorded orders, none of the three 1,024-draw menus cost an incumbent their placement. In two cohorts the old assignment nevertheless lost its stability certificate. The eight edit-matched comparison seeds produced 894--968 gross exits, summing three cohorts separately for each seed. Table~\ref{tab:nested-risk} uses the same simulation holdout we had already inspected at 16 draws.

We then froze the 1,024-draw design and ran 256 additional draws per cohort with a fresh seed batch. Across these 768 cohort--lottery realizations the screen caused 87 gross exits and 91 entrants. There was at least one exit in 26, 34 and 23 draws by year, and at most two exits in any single draw. The certificate failed in 55, 76 and 56 draws. The edit-matched control using the first seed caused 237,115 gross exits across the same 768 realizations. At the recorded orders, the screened menus improved net placement by one, placed five baseline recipients at worse ranks and gave two fewer first choices. They showed 14 more initial offers and 27 fewer April schools than the blind predictor. The unrestricted switch showed 4,228 more initial offers and 3,590 more April schools. Each historical label counts an applicant at most once, including where more than one listed course maps to the same school.

\paragraph{Could another protected menu recover the predictive gain?}
For each historical target separately, imagine knowing every pupil's label. If the target already appears on $A_i$, keep that row. Otherwise one new hit is possible exactly when (i) some proposed-only option matches that target and passes every protected draw, and (ii) at least one old-only option is removable. These conditions use fixed assignments and hold independently by row, so summing the indicators gives the exact maximum increase over the blind menu. A course-level initial offer and a school-level April enrollment each get their own maximizing display. This is an information upper bound over the union of the two menus with shared choices protected.

With no lottery constraint, the union permitted up to 4,823 additional April-school hits. Requiring stability in all 16, 64, 256 and 1,024 design draws lowered that number to 593, 353, 302 and 278 respectively. For the last support, the maximum initial-offer gain was 48. Even the best feasible use of the historical labels could reproduce at most $278/3{,}590=7.7\%$ of the unrestricted rule's net increase in April-school coverage. This limit concerns the stated menu pool and the requirement to preserve each old assignment as a stable outcome. Protecting only whether a pupil remains placed would give a different design problem.
